\section{Tài liệu tham khảo}

{[}1{]} Khoa Công nghệ Thông tin, Trường ĐH Khoa học Tự nhiên -- ĐHQG-HCM, \emph{CSC14003 Project 1: Bloxorz Solver} (đề bài), 2026.

{[}2{]} D. Clarke, \emph{Bloxorz} (trò chơi Flash), 2007. Coolmath Games. {[}Trực tuyến{]}. Địa chỉ: \url{https://www.coolmathgames.com/0-bloxorz}

{[}3{]} J. Coughenour, \emph{bloxorz\_gba} (bản port Game Boy Advance của Bloxorz, 33 màn gốc), GitHub, commit f57ebb6, 2026, giấy phép GPL-3.0. {[}Trực tuyến{]}. Địa chỉ: \url{https://github.com/jacobcoughenour/bloxorz_gba}

{[}4{]} J. Gordon, \emph{quick\_bloxorz}, GitHub, commit 07821f5, 2021, giấy phép CC BY. {[}Trực tuyến{]}. Địa chỉ: \url{https://github.com/jgordon510/quick_bloxorz}

{[}5{]} O. Cheong, \emph{cs109: Bloxorz levels}, GitHub, commit 0097f4a, 2015. {[}Trực tuyến{]}. Địa chỉ: \url{https://github.com/otfried/cs109/tree/master/bloxorz/levels}

{[}6{]} L. Pham và cộng sự, \emph{BloxorzSolver}, GitHub, commit cd9c142, 2018. {[}Trực tuyến{]}. Địa chỉ: \url{https://github.com/phamquiluan/BloxorzSolver}

{[}7{]} S. Russell và P. Norvig, \emph{Artificial Intelligence: A Modern Approach}, 4th ed.~Hoboken, NJ, USA: Pearson, 2020.

{[}8{]} P. T. Tran và cộng sự, \emph{bloxorz-game} (đồ án Bloxorz Solver, môn Cơ sở Trí tuệ nhân tạo, ĐH Khoa học Tự nhiên -- ĐHQG-HCM), GitHub, commit 4f3afeb, 2026. {[}Trực tuyến{]}. Địa chỉ: \url{https://github.com/TranTai35/bloxorz-game}

{[}9{]} Python Software Foundation, \emph{Python 3.11}. {[}Trực tuyến{]}. Địa chỉ: \url{https://www.python.org}

{[}10{]} P. Säde, \emph{Ursina Engine}, phiên bản 7.0.0. {[}Trực tuyến{]}. Địa chỉ: \url{https://www.ursinaengine.org}

{[}11{]} Carnegie Mellon University và Walt Disney Imagineering, \emph{Panda3D}, phiên bản 1.10.16. {[}Trực tuyến{]}. Địa chỉ: \url{https://www.panda3d.org}

{[}12{]} J. D. Hunter, ``Matplotlib: A 2D graphics environment,'' \emph{Computing in Science \& Engineering}, vol.~9, no. 3, pp.~90--95, 2007. (Dùng phiên bản 3.10.9.)

{[}13{]} The Be Vietnam Pro Project Authors, \emph{Be Vietnam Pro} (phông chữ của giao diện), giấy phép SIL Open Font License 1.1. {[}Trực tuyến{]}. Địa chỉ: \url{https://github.com/bettergui/BeVietnamPro}

{[}14{]} three.js authors, \emph{three.js} r160, kèm \texttt{OrbitControls.\allowbreak{}js} và \texttt{RoundedBoxGeometry.\allowbreak{}js}, giấy phép MIT. {[}Trực tuyến{]}. Địa chỉ: \url{https://threejs.org}

{[}15{]} The Baloo 2 Project Authors, \emph{Baloo 2}, và The Nunito Project Authors, \emph{Nunito} (phông chữ), giấy phép SIL Open Font License 1.1. {[}Trực tuyến{]}. Địa chỉ: \url{https://github.com/EkType/Baloo2}, \url{https://github.com/googlefonts/nunito}

\textbf{Ghi chú về nguồn 8.} Bài của một nhóm khóa trước cùng môn, không có giấy phép. Nhóm chỉ lấy ý tưởng (bài học về thực nghiệm, nút So sánh, các lỗi hay gặp với Ursina, sân trống lớn và màn không lời giải), không dùng code, asset hay dữ liệu màn.

\textbf{Ghi chú về nguồn 3 đến 6.} Dữ liệu 33 màn gốc được chuyển tự động từ nguồn {[}3{]} (giấy phép GPL-3.0); ba nguồn {[}4{]}, {[}5{]}, {[}6{]} chỉ dùng để đối chiếu từng ô của các màn (mục 5.2). \textbf{Nguồn 14 và 15} chỉ dùng trong bản demo giao diện web đầu tiên (thư mục \texttt{Source/\allowbreak{}web/\allowbreak{}}), không dùng trong giao diện Ursina.
